\paragraph{From the incidence origin to lattice modification.}
\label{inc-ampl-reticulo}
The marked flag determines a position through an equivariant
set-theoretic operation:
\[
 o_*=(H_e\cap H_\varphi)
       \setminus(H^-_\alpha\cup H_{\rm APP})=\{1\}.
\]
Its function is to distinguish one component among the twelve
involved in the lattice realization. The complete ternary code
acts as a glue code in the discriminant module
$(A_2^*/A_2)^{12}$. Its preimage defines $N(C_W)$, and
the self-duality, weights, and distance of the code translate,
respectively, into integrality and determinant properties, parity,
and norm bounds for the glue classes. Here the symbolic information
of the code serves a function that the isolated support no longer performs.

The $A_2$ metric and the positive radial chamber are fixed in the
lattice development. Within that chamber, the congruence required for an
even index-three neighbour has twelve realizations of minimum norm $54$,
permuted by changing the distinguished component. The origin $o_*$
selects
\[
 v_\alpha=4\rho_1+\rho_2+\cdots+\rho_{12},\qquad
 v_\alpha^2=2(4^2+11)=54.
\]
The coefficient $4$ belongs to that solution of the congruence in the
stated chamber. The subscript $\alpha$ records the incidence
provenance of the selection: its content is the frame that participated in
the closure and now singles out a lattice direction.

Pairing with $v_\alpha$ determines the index-three sublattice
$N_{\alpha,0}$, and adjoining the class
$v_\alpha/3$ produces the neighbour $N_\alpha$ of
\eqref{exc:vecino}. This modification preserves rank and recovers
unimodularity and evenness; congruence selection eliminates the
previous roots, and local bounds exclude roots in the new classes.
Identification with Leech follows these verifications. The
binary realization of the five regions and this lattice construction
are two prolongations of the common frame: the latter directly receives
the ternary code as glue data.
